%% 
%% Copyright 2019-2024 Elsevier Ltd
%% 
%% Template article for cas-dc documentclass for double column output.
%% 

\documentclass[a4paper,fleqn]{cas-sc}

\usepackage{amsmath,amsfonts,amssymb,amsthm}
\usepackage{algorithmic}
\usepackage{algorithm}
\usepackage{array}
\usepackage{textcomp}
\usepackage{stfloats,siunitx}
\usepackage{url}
\usepackage{verbatim}
\usepackage{graphicx}
\graphicspath{{figs/}}
\usepackage{cite,soul}
\usepackage{multirow,booktabs}
\usepackage{float,booktabs,threeparttable} 
\usepackage{tikz,graphics,color,fullpage,float,epsf,caption,subcaption}
\usepackage{epsfig,epstopdf,arydshln,microtype}

\usepackage[numbers]{natbib}

%%% Author macros
\def\tsc#1{\csdef{#1}{\textsc{\lowercase{#1}}\xspace}}
\tsc{WGM}
\tsc{QE}

\newtheorem{theorem}{Theorem}
\newtheorem{lemma}[theorem]{Lemma}
\newtheorem{corollary}{Corollary}
\newdefinition{rmk}{Remark}
\newproof{pf}{Proof}

\begin{document}
\let\WriteBookmarks\relax
\def\floatpagepagefraction{1}
\def\textpagefraction{.001}

% Short title
\shorttitle{Resource-Bounded Non-Parametric SPC Framework for High-Dimensional Non-Gaussian IoT Stream}
%\shorttitle{Resource-Bounded Non-Parametric Bootstrap Control Chart via RBULT}    


% Short author
\shortauthors{P. Junsawang et~al.}  

% Main title of the paper
\title [mode = title]{Resource-Bounded Non-Parametric SPC Framework for High-Dimensional Non-Gaussian IoT Stream}
%\title [mode = title]{Resource-Bounded Non-Parametric Bootstrap Control Chart for High-Dimensional Non-Gaussian Data Streams via RBULT}  

\tnotemark[1] 
\tnotetext[1]{This work was supported in part by the Advanced Virtual and Intelligent Computing (AVIC) Research Center, Department of Mathematics and Computer Science, Faculty of Science, Chulalongkorn University.} 

% Authors
\author[1]{P. Junsawang}
\ead{jpremlc@gmail.com}

\author[1]{S. Phimoltares}
\cormark[1]
\cortext[1]{Corresponding author}
\ead{suphakant.p@chula.ac.th}

\author[1]{C. Lursinsap}
\ead{chidchanok.L@chula.ac.th}

\affiliation[1]{organization={Advanced Virtual and Intelligent Computing (AVIC) Research Center, Department of Mathematics and Computer Science, Faculty of Science, Chulalongkorn University},
            city={Bangkok},
            postcode={10330}, 
            country={Thailand}}

% Abstract
\begin{abstract}
Bootstrapping is a powerful non-parametric technique for quantifying uncertainty without restrictive distributional assumptions. However, applying conventional bootstrap methods to real-time industrial data streams poses two fundamental challenges: prohibitive memory overhead ($O(N \cdot D)$ storage growth) and susceptibility to control limit pollution when anomalous samples contaminate historical memory buffers. This paper presents \textbf{RBULT-SPC}, a resource-bounded non-parametric control chart framework designed for high-dimensional, non-Gaussian IoT data streams. Grounded in mathematical proofs of chunk extreme probabilities, RBULT-SPC maintains only compact tail summary statistics in $O(D)$ constant space, preventing memory overflow and eliminating buffer pollution. To handle high dimensionality ($D=34$), the framework combines Z-score outlier filtering, lazy boundary expansion, 11-candidate non-parametric distribution MLE fitting, and Bonferroni Family-Wise Error Rate (FWER) tail adjustment. Extensive two-tier empirical evaluations across 10 simulated 1D distributions, 4 economic datasets, 5 industrial streams, and 4 operating regimes of the Tennessee Eastman Process (TEP) demonstrate that RBULT-SPC achieves optimal coverage ($93.71\% - 99.95\%$) and controlled false alarm rates ($0.05\% - 6.29\%$), outperforming classical Shewhart and EWMA charts whose false alarm rates collapse to $18.21\% - 39.01\%$ under non-Gaussianity and operational throughput stress. Furthermore, hyperparameter sensitivity analysis demonstrates that dimension-aware thresholding completely eliminates batch false alarms ($\text{Chunk FAR} = 0.00\%$) while maintaining an immediate failure detection response delay ($\text{ARL}_1 = 1.00$) at an average latency under 45 ms.
\end{abstract}

% Research highlights
\begin{highlights}
\item We propose \textbf{RBULT-SPC}, a novel resource-bounded non-parametric control chart framework for high-dimensional non-Gaussian data streams.
\item Combines foundational 1D probabilistic chunk excursion proofs with multivariate Bonferroni Family-Wise Error Rate (FWER) tail scaling.
\item Achieves strict $O(D)$ constant RAM memory boundedness ($3.23\text{ KB}$), delivering a $180\times$ memory reduction over conventional sliding-window bootstrap ($582.87\text{ KB}$).
\item Integrates Algorithm 4 Z-score spike filtering with dynamic 11-candidate distribution MLE fitting to prevent control limit pollution.
\item Eliminates batch false-alarm spam ($\text{Chunk FAR} = 0.00\%$) on high-dimensional streams ($D=34$) while retaining an immediate anomaly detection response delay ($\text{ARL}_1 = 1.00$).
\item Validated across a comprehensive two-tier benchmark: 10 simulated 1D distributions, 4 economic datasets, 5 real industrial streams, and 4 TEP stress regimes (Modes 1, 3, 4, 5) up to 1.74M observations.
\end{highlights}

% Keywords
\begin{keywords}
Online Bootstrapping \sep Statistical Process Control \sep High-Dimensional IoT Data Streams \sep Non-Gaussian Resampling \sep Limited Memory Environments
\end{keywords}

\maketitle

% =========================================================================
% SECTION 1: INTRODUCTION
% =========================================================================
\section{Introduction}\label{sec1_intro}
In modern smart factory environments and Industry 4.0 applications, real-time quality control relies heavily on monitoring continuous sensor telemetry (e.g., pressure, temperature, current, vibration, rotational speed). Statistical Process Control (SPC) charts serve as the foundational mechanism for detecting out-of-control (OOC) states and operational faults.

\subsection{Industrial Context \& Problem Motivation}
Classical SPC techniques, including Shewhart $\bar{X}$-charts, EWMA, and CUSUM, rely on the parametric assumption that underlying telemetry streams follow a Gaussian normal distribution. However, real-world sensor streams in industrial processes frequently exhibit non-Gaussian characteristics, including extreme skewness, heavy tails, and negative kurtosis.

\subsection{Critical Bottlenecks of Existing Methods}
Applying traditional resampling and bootstrap methods to high-dimensional streaming environments encounters three primary bottlenecks:
\begin{enumerate}
    \item \textbf{Memory Storage Explosion ($O(N \cdot D)$):} Storing historical observations indefinitely leads to Out-Of-Memory (OOM) failures on resource-constrained embedded IoT edge devices.
    \item \textbf{Control Limit Pollution:} When anomalous fault observations contaminate sliding memory buffers, conventional bootstrap limits expand unnaturally, masking subsequent anomalies.
    \item \textbf{High-Dimensional False Alarm Spam:} Monitoring $D$ sensor channels simultaneously without Family-Wise Error Rate (FWER) adjustment results in an unacceptable rate of false alarms.
\end{enumerate}

\subsection{Summary of Contributions}
To address these challenges, this paper presents \textbf{RBULT-SPC}. The major contributions of this work include:
\begin{itemize}
    \item Mathematical formulation and proofs of univariate chunk extreme occurrence probabilities (Theorems 1--2, Corollary 1) and multivariate FWER control (Theorem 3).
    \item A strict $O(D)$ constant RAM architecture ($3.23\text{ KB}$), achieving a $180\times$ RAM reduction over sliding-window bootstrap ($582.87\text{ KB}$).
    \item A 4-stage statistical pipeline integrating Z-score spike filtering, lazy boundary expansion, 11-candidate distribution MLE fitting, and Bonferroni tail adjustment.
    \item Extensive two-tier empirical validation across 10 synthetic 1D distributions, 4 economic datasets, 5 industrial telemetry streams, and 4 TEP operating regimes ($D=34$, up to 1.74M observations).
\end{itemize}


% =========================================================================
% SECTION 2: RELATED WORK
% =========================================================================
\section{Related Work}\label{sec2_related}
This section reviews relevant literature across three domain areas:
\subsection{Classical and Multivariate Statistical Process Control}
Review of Shewhart, EWMA, CUSUM, and Hotelling's $T^2$ control charts, highlighting performance degradation under non-Gaussianity.

\subsection{Bootstrap Resampling in Streaming Environments}
Review of Efron's percentile bootstrap, blockwise bootstrap, and online streaming bootstrap implementations, noting memory overhead and buffer pollution risks.

\subsection{Uncertainty Quantification and Outlier Filtering in IoT Data Streams}
Review of streaming range estimation, hyper-rectangle bounding boxes, and Z-score outlier filtering techniques.


% =========================================================================
% SECTION 3: THEORETICAL FOUNDATIONS & 1D RANGE ESTIMATION
% =========================================================================
\section{Theoretical Foundations and Univariate Range Estimation}\label{sec3_theory}
This section presents the mathematical foundation and 1D range estimation formulation of the proposed Recursively Bootstrapping at Upper-Lower Tails (RBULT) concept, establishing the theoretical ground truth for streaming boundary expansion.

\subsection{Concerned Scenario}
\label{scenario}
Sampled data continuously flow in forms of chunks into a process, which may be a learning process or a statistical analyzing process. The size of each chunk is not fixed and does not exceed the size of the working memory. Each sampled chunk contains a set of positive numbers whose minimum and maximum values are unknown. The probability distribution of each integer or continuous value is not predefined and unknown. This scenario exists in various problems in machine learning and data stream analytics. Estimating the population range or variance of incoming data in advance under the \emph{discard-after-learn concept} helps establish adaptive boundaries without retaining past raw observations.

\subsection{Studied Problems}
Suppose the size of working memory is fixed at $m$. Let $c_{i}$ be the $i^{th}$ incoming numeric data chunk such that $|c_{i}| < m$ and $i \geq 1$. After receiving chunk $c_{i}$ at time $t = i$, the minimum and maximum values of population at the left and right ends are estimated. Let $L$ denote the estimated left-end value which may be the minimum data value of chunk $c_{i}$ or the minimum value computed from bootstrapping data in $c_{i}$, depending on which one is smaller. Similar to $L$, let $R$ denote the estimated right-end value which may be the maximum data value of chunk $c_{i}$ or the maximum value computed from bootstrapping data in $c_{i}$, depending on which one is larger. After estimating $L$ and $R$ of $c_{i}$, all data except partial summary statistics at both ends are discarded.

\begin{enumerate}
    \item How to estimate the population variance interval without using the entire data in working memory set $D$ so that data in subsequent incoming chunks fall within the new estimated interval $[L, R]$?
    \item Theoretically, how many times on average must the data interval in set $D$ be expanded according to the estimated population variance?
\end{enumerate}

\subsection{Constraints}
Let $|c_i|$ represent the number of elements in a given chunk $c_i$. The bin set $B$ is an interval with a left end $L$ and a right end $R$. The following constraints apply:
\begin{enumerate}
    \item The probability distribution of data within a chunk $c_i$ is unknown.
    \item The size of chunk $|c_i|$ may vary across streaming iterations, with $|c_{i}| < |D|$.
    \item Once all data in chunk $c_i$ have been evaluated, the raw data chunk is completely discarded and cannot be re-inspected.
\end{enumerate}

\subsection{Concept of Recursively Bootstrapping at Upper-Lower Tails (RBULT)}
\label{concept}
RBULT resamples data only in the left-end and right-end tail intervals ($b_1$ and $b_8$), differing from classical bootstrap which resamples across the entire data body. The size of left-end and right-end intervals is determined from the standard histogram shape of the best-fitting probability distribution. If the heights of left-end ($|b_1|$) or right-end ($|b_8|$) intervals in the actual incoming data histogram exceed the theoretical heights ($|h_1|$ or $|h_8|$) of the standard histogram, the boundary intervals are extended by bootstrapping data in $b_1$ and $b_8$.

The RBULT execution operates in two phases:
\begin{enumerate}
    \item \textbf{Phase 1 (Initial Fitting):} Analyze the probability distribution of the first chunk $c_1$ and compare the standard histogram with the actual histogram. If heights exceed standard thresholds, bootstrap $b_1$ and $b_8$ to extend bounds $[L, R]$.
    \item \textbf{Phase 2 (Streaming Recurrence):} For each new incoming chunk $c_i$ ($i > 1$), update interval bins, evaluate boundary excursions, and recursively expand $[L, R]$ whenever actual tail heights exceed standard thresholds. Discard non-tail data afterwards.
\end{enumerate}

\subsection{Candidate Distribution Selection and CDF Tail Bin Area Integration}
To capture diverse physical telemetry behaviors in smart manufacturing (e.g., right-skewed bearing friction, Weibull mechanical wear, Wald diffusion processes, and negative-kurtosis uniform signals), RBULT evaluates a candidate family of continuous distributions $\mathcal{P} = \{\text{Gamma}$, $\text{ExponWeibull}$, $\text{Weibull}_{\min/\max}$, $\text{Wald}$, $\text{ExponPow}$, $\text{Rayleigh}$, $\text{Powerlaw}$, $\text{Lognormal}$, $\text{Chi-Square}$, $\text{Uniform}$, $\text{Normal}\}$. Each candidate represents distinct tail skewness and kurtosis characteristics observed in real-world IoT sensor streams.
To construct the standard histogram $H = \{h_1, \dots, h_8\}$ without relying on standard normal assumptions, the candidate family of continuous distributions is evaluated on initial data chunk $c_1$. The candidate distribution achieving the optimal parameter fit $\hat{\boldsymbol{\theta}}$ is then selected by minimizing the Akaike Information Criterion (AIC), defined as:
\begin{equation}
\text{AIC} = 2k - 2\ln \mathcal{L}(\hat{\boldsymbol{\theta}} \mid c_1)
\end{equation}
where $k$ represents the number of estimated parameters (complexity penalty) and $\ln \mathcal{L}(\hat{\boldsymbol{\theta}} \mid c_1)$ denotes the maximized log-likelihood value.



For each candidate density $f(x; \boldsymbol{\theta}) \in \mathcal{P}$, the parameter vector $\boldsymbol{\theta}$ is estimated via Maximum Likelihood Estimation (MLE) by maximizing the log-likelihood function:
\begin{equation}
\hat{\boldsymbol{\theta}} = \arg\max_{\boldsymbol{\theta}} \ln \mathcal{L}(\boldsymbol{\theta} \mid c_1) = \arg\max_{\boldsymbol{\theta}} \sum_{j=1}^{|c_1|} \ln f(x_{1,j}; \boldsymbol{\theta})
\end{equation}
The candidate distribution achieving the optimal parameter fit $\hat{\boldsymbol{\theta}}$ is then selected by minimizing the Akaike Information Criterion (AIC), defined as:
\begin{equation}
\text{AIC} = 2k - 2\ln \mathcal{L}(\hat{\boldsymbol{\theta}} \mid c_1)
\end{equation}
where $k$ represents the number of estimated parameters (complexity penalty) and $\ln \mathcal{L}(\hat{\boldsymbol{\theta}} \mid c_1)$ denotes the maximized log-likelihood value. Minimizing AIC ($\arg\min \text{AIC}$) selects the optimal candidate distribution $f(x; \hat{\boldsymbol{\theta}}) \in \mathcal{P}$ that achieves the best trade-off between fitting accuracy and model simplicity (parsimony). This prevents selecting unnecessarily complex distributions that might overfit transient noise spikes in initial data chunk $c_1$. Once $\hat{\boldsymbol{\theta}}$ is selected, the expected theoretical element count $|h_k|$ for each bin $b_k$ ($k \in \{1, \dots, 8\}$) is computed via exact CDF integration:

\begin{equation}
|h_1| = |c_1| \cdot \int_{\mu - 4\sigma}^{\mu - 3\sigma} f(x; \hat{\boldsymbol{\theta}}) \, dx = |c_1| \left[ F(\mu - 3\sigma; \hat{\boldsymbol{\theta}}) - F(\mu - 4\sigma; \hat{\boldsymbol{\theta}}) \right]
\end{equation}
\begin{equation}
|h_8| = |c_1| \cdot \int_{\mu + 3\sigma}^{\mu + 4\sigma} f(x; \hat{\boldsymbol{\theta}}) \, dx = |c_1| \left[ F(\mu + 4\sigma; \hat{\boldsymbol{\theta}}) - F(\mu + 3\sigma; \hat{\boldsymbol{\theta}}) \right]
\end{equation}

If the actual observed data counts $|b_1|$ or $|b_8|$ in chunk $c_1$ exceed $|h_1|$ or $|h_8|$, the boundary expansion loop (Algorithm 3) is triggered.

\subsection{RBULT Algorithms}
This subsection details the 4 algorithmic specifications governing 1D RBULT streaming boundary estimation, dual-protocol tracking, and Z-score spike filtering.

\begin{algorithm}[H]
\caption{Initial Chunk Fitting and Tail Summary Setup ($c_1$)}
\label{alg1}
\begin{algorithmic}[1]
\REQUIRE First data chunk $c_1$, List $\mathcal{P}$ of candidate probability functions, Tail interval set $B=\{b_1, \dots, b_8\}$, Standard histogram $H=\{h_1, \dots, h_8\}$.
\ENSURE Initial left-end bound $L_1$ and right-end bound $R_1$.
\STATE Load $c_1$ into working buffer $D$.
\STATE Estimate parameters $\hat{\boldsymbol{\theta}}$ via MLE and select best distribution $f(x; \hat{\boldsymbol{\theta}}) \in \mathcal{P}$ minimizing AIC/KS.
\STATE Compute expected tail counts $|h_1|$ and $|h_8|$ via exact CDF integration over $[\mu - 4\sigma, \mu - 3\sigma]$ and $[\mu + 3\sigma, \mu + 4\sigma]$.
\STATE Initialize $L_1 = \min(D)$ and $R_1 = \max(D)$.
\WHILE{$|b_1| > |h_1|$ \textbf{or} $|b_8| > |h_8|$}
    \STATE Bootstrap tail intervals $b_1$ and $b_8$ via Algorithm \ref{alg3} to refine $L_1$ and $R_1$.
    \STATE Update bin center $\mu = \frac{L_1 + R_1}{2}$ and width $\sigma = \frac{R_1 - L_1}{8}$, and recompute $B$ and $H$.
\ENDWHILE
\STATE Store tail summary statistics in $b_1, b_8$ and discard raw chunk data from $D$.
\end{algorithmic}
\end{algorithm}

\begin{algorithm}[H]
\caption{Streaming Recurrent Boundary Update with Dual Protocol Tracking ($c_m, m > 1$)}
\label{alg2}
\begin{algorithmic}[1]
\REQUIRE Streaming chunk $c_m$, Existing bounds $L_{m-1}, R_{m-1}$, Tail set $B$, Outlier threshold $\tau$.
\ENSURE Updated bounds $L_m, R_m$, Pre-Sequential bounds $[L_m^{\text{pre}}, R_m^{\text{pre}}]$, In-Sample bounds $[L_m^{\text{in}}, R_m^{\text{in}}]$.
\WHILE{new streaming chunk $c_m$ arrives}
    \STATE Load $c_m$ into working buffer $D$.
    \STATE \textbf{Record Pre-Sequential Bounds (Pre-Update Predictive Bounds):}
    \STATE \quad Set $[L_m^{\text{pre}}, R_m^{\text{pre}}] = [L_{m-1}, R_{m-1}]$.
    \STATE Filter transient outlier spikes from $D$ via Algorithm \ref{alg4} using threshold $\tau$.
    \STATE \textbf{Lazy Boundary Expansion:}
    \IF{$\min(D) < L_{m-1}$} \STATE $L_m = \min(D)$. \ELSE \STATE $L_m = L_{m-1}$. \ENDIF
    \IF{$\max(D) > R_{m-1}$} \STATE $R_m = \max(D)$. \ELSE \STATE $R_m = R_{m-1}$. \ENDIF
    \STATE Reconstruct data histogram $B$ and standard histogram $H$ using $\mu = \frac{L_m + R_m}{2}$ and $\sigma = \frac{R_m - L_m}{8}$.
    \WHILE{$|b_1| > |h_1|$ \textbf{or} $|b_8| > |h_8|$}
        \STATE Resample tail bins $b_1$ and $b_8$ to expand $L_m$ and $R_m$ via Algorithm \ref{alg3}.
        \STATE Recompute $\mu, \sigma$ and update histograms $B$ and $H$.
    \ENDWHILE
    \STATE \textbf{Record In-Sample Bounds (Post-Update Adapted Bounds):}
    \STATE \quad Set $[L_m^{\text{in}}, R_m^{\text{in}}] = [L_m, R_m]$.
    \STATE Discard raw stream samples $D$ to preserve $O(1)$ constant RAM storage.
\ENDWHILE
\end{algorithmic}
\end{algorithm}

\begin{algorithm}[H]
\caption{Recurrent Tail Resampling for Boundaries $L$ and $R$}
\label{alg3}
\begin{algorithmic}[1]
\REQUIRE Tail data intervals $b_1$ and $b_8$, Bootstrap sample size $N_{\text{boot}}$, Iterations $M_{\text{iter}}$.
\ENSURE Refined left bound $L$ and right bound $R$.
\STATE \textbf{--- Left Tail Boundary Expansion ($L$ via $b_1$) ---}
\STATE Initialize mean set $\mathcal{M}_{b_1} = \emptyset$ and std set $\mathcal{S}_{b_1} = \emptyset$.
\FOR{$k = 1$ \TO $M_{\text{iter}}$}
    \STATE Sample $N_{\text{boot}}$ points from $b_1$ with replacement; calculate sample mean $\hat{\mu}_{k,b_1}$ and std $\hat{\sigma}_{k,b_1}$.
    \STATE Update $\mathcal{M}_{b_1} = \mathcal{M}_{b_1} \cup \{\hat{\mu}_{k,b_1}\}$ and $\mathcal{S}_{b_1} = \mathcal{S}_{b_1} \cup \{\hat{\sigma}_{k,b_1}\}$.
\ENDFOR
\STATE Compute bootstrap mean $\mu_{b_1}^{(\text{boot})} = \text{mean}(\mathcal{M}_{b_1})$ and dispersion $\sigma_{b_1}^{(\text{boot})} = \text{mean}(\mathcal{S}_{b_1})$.
\IF{$\text{mean}(b_1) < \mu_{b_1}^{(\text{boot})}$}
    \STATE $L = \min(b_1) - |\text{std}(b_1) - \sigma_{b_1}^{(\text{boot})}|$.
\ELSE
    \STATE $L = \min(b_1) - |\text{mean}(b_1) - \mu_{b_1}^{(\text{boot})}|$.
\ENDIF
\STATE \textbf{--- Right Tail Boundary Expansion ($R$ via $b_8$) ---}
\STATE Initialize mean set $\mathcal{M}_{b_8} = \emptyset$ and std set $\mathcal{S}_{b_8} = \emptyset$.
\FOR{$k = 1$ \TO $M_{\text{iter}}$}
    \STATE Sample $N_{\text{boot}}$ points from $b_8$ with replacement; calculate sample mean $\hat{\mu}_{k,b_8}$ and std $\hat{\sigma}_{k,b_8}$.
    \STATE Update $\mathcal{M}_{b_8} = \mathcal{M}_{b_8} \cup \{\hat{\mu}_{k,b_8}\}$ and $\mathcal{S}_{b_8} = \mathcal{S}_{b_8} \cup \{\hat{\sigma}_{k,b_8}\}$.
\ENDFOR
\STATE Compute bootstrap mean $\mu_{b_8}^{(\text{boot})} = \text{mean}(\mathcal{M}_{b_8})$ and dispersion $\sigma_{b_8}^{(\text{boot})} = \text{mean}(\mathcal{S}_{b_8})$.
\IF{$\text{mean}(b_8) > \mu_{b_8}^{(\text{boot})}$}
    \STATE $R = \max(b_8) + |\text{std}(b_8) - \sigma_{b_8}^{(\text{boot})}|$.
\ELSE
    \STATE $R = \max(b_8) + |\text{mean}(b_8) - \mu_{b_8}^{(\text{boot})}|$.
\ENDIF
\end{algorithmic}
\end{algorithm}

\begin{algorithm}[H]
\caption{Z-Score Outlier Spike Filtering for Contaminated Streams}
\label{alg4}
\begin{algorithmic}[1]
\REQUIRE Data buffer $D$, Spike sensitivity threshold $\tau$ (default $\tau = 3.0$).
\ENSURE Filtered data buffer $D$ free of transient noise spikes.
\STATE Compute sample mean $\mu_D = \frac{1}{|D|}\sum_{x \in D} x$ and sample std $\sigma_D = \sqrt{\frac{1}{|D|-1}\sum_{x \in D}(x - \mu_D)^2}$.
\FOR{each sample $x \in D$}
    \STATE Calculate standard score $z(x) = \frac{|x - \mu_D|}{\sigma_D}$.
    \IF{$z(x) > \tau$}
        \STATE Exclude outlier spike $x$ from buffer $D$.
    \ENDIF
\ENDFOR
\STATE Return cleaned working buffer $D$.
\end{algorithmic}
\end{algorithm}

\subsection{Theoretical Analysis}
The theoretical minimum and maximum numbers of boundary expansion can be analyzed in terms of probability combinatorics. Let population size be $N$ and chunk size be $l$ without duplicate observations. Let $p^{(\min)} = \min(\mathcal{P})$ and $p^{(\max)} = \max(\mathcal{P})$.

\begin{theorem}[First Chunk Extreme Probability]
Assume a chunk of $l$ observations is sampled without replacement from population $\mathcal{P}$ of size $N$. The probability of containing both $p^{(\min)}$ and $p^{(\max)}$ in the first chunk is:
\begin{equation}
P(E_1) = \frac{l(l-1)}{N(N-1)}
\end{equation}
\end{theorem}
\begin{proof}
When $p^{(\min)}$ and $p^{(\max)}$ are present in the first chunk, the remaining $l-2$ observations must be chosen from the remaining $N-2$ population elements. The number of ways to pick $l-2$ elements from $N-2$ is $\binom{N-2}{l-2}$. Combining with $(N-l)!$ permutations of remaining population items yields $c = \binom{N-2}{l-2}(N-l)! = \frac{(N-2)!}{(l-2)!}$. Incorporating $l!$ chunk permutations yields $l(l-1)(N-2)!$ total favorable permutations out of $N!$ possible population arrangements. Thus:
\begin{equation}
P(E_1) = \frac{l(l-1)(N-2)!}{N!} = \frac{l(l-1)}{N(N-1)}
\end{equation}
\end{proof}

\begin{theorem}[Streaming Chunk Distribution]
Assume a chunk of $l$ observations is sampled without replacement from population $\mathcal{P}$ of size $N$. The probability of having $p^{(\min)}$ and $p^{(\max)}$ in the $i^{th}$ incoming chunk ($i \ge 1$) is uniform and equal to:
\begin{equation}
P(E_i) = \frac{l(l-1)}{N(N-1)}
\end{equation}
\end{theorem}
\begin{proof}
Let $S$ denote the sample space of partitioning $N$ observations into $L$ chunks of size $l$ ($N = l \cdot L$). The total number of sample outcomes is $|S| = \frac{N!}{(l!)^L}$. Let $E_i \subset S$ denote events where both extreme bounds $p^{(\min)}$ and $p^{(\max)}$ reside in chunk $i$. By product rule combinatorics:
\begin{equation}
|E_i| = \binom{N-2}{l-2} \frac{(N-l)!}{(l!)^{L-1}}
\end{equation}
Evaluating the probability ratio yields:
\begin{equation}
P(E_i) = \frac{|E_i|}{|S|} = \frac{\binom{N-2}{l-2}\frac{(N-l)!}{(l!)^{L-1}}}{\frac{N!}{(l!)^L}} = \frac{l(l-1)}{N(N-1)}
\end{equation}
which proves uniform occurrence probability across all streaming chunk positions $i \ge 1$.
\end{proof}

\begin{corollary}[Expected Chunk Index for Extreme Bounds]
Assuming joint occurrence of minimum and maximum observations is restricted to a single chunk, the expected chunk index $\mathbb{E}[X]$ containing both extreme bounds across $L$ total chunks is:
\begin{equation}
\mathbb{E}[X] = \frac{1}{2} \left( \frac{l(l-1) L (L+1)}{N(N-1)} \right)
\end{equation}
\end{corollary}
\begin{proof}
Let $X \in \{1, 2, \dots, L\}$ be a discrete random variable representing the chunk index containing $p^{(\min)}$ and $p^{(\max)}$. Summing over uniform probability $P(E_i) = \frac{l(l-1)}{N(N-1)}$ yields:
\begin{align}
\mathbb{E}[X] &= \sum_{i=1}^{L} i \cdot P(X = i) = \left( \sum_{i=1}^{L} i \right) \frac{l(l-1)}{N(N-1)} \notag \\
&= \left( \frac{L(L+1)}{2} \right) \left( \frac{l(l-1)}{N(N-1)} \right) = \frac{1}{2} \left( \frac{l(l-1) L (L+1)}{N(N-1)} \right)
\end{align}
\end{proof}


% =========================================================================
% SECTION 4: PROPOSED HIGH-DIMENSIONAL RBULT-SPC FRAMEWORK
% =========================================================================
\section{Proposed High-Dimensional RBULT-SPC Framework}\label{sec4_framework}
This section presents the extension of the 1D RBULT concept to multivariate high-dimensional Statistical Process Control ($\text{RBULT-SPC}$).

\subsection{System Architecture and 4-Stage Streaming Pipeline}
Details of the 4-stage streaming process: Stationary Preprocessing, Z-score Spike Filter, Dimensional Bound Estimation, and Bounding Hyper-Rectangle $\mathcal{B}_m = \prod_{d=1}^D [L_{m,d}, R_{m,d}]$ evaluation.

\subsection{Bonferroni / Šidák Family-Wise Error Rate (FWER) Scaling}
To maintain a target system false alarm rate $\alpha_{\text{sys}} = 0.05$ across $D$ monitored sensor channels, per-dimension tail quantiles are scaled using Bonferroni correction:
\begin{equation}
\alpha_{\text{dim}} = \frac{\alpha_{\text{sys}}}{D}
\end{equation}
yielding dimension-aware control limit formulations:
\begin{equation}
\text{LCL}_d = L_d = Q_{\text{tail}}\left(\frac{\alpha_{\text{sys}}}{2D}\right), \quad \text{UCL}_d = R_d = Q_{\text{tail}}\left(1 - \frac{\alpha_{\text{sys}}}{2D}\right)
\end{equation}

\subsection{Theorem 3 (Family-Wise Error Rate Upper Bound)}
\begin{theorem}[Multivariate FWER Upper Bound]
For a $D$-dimensional streaming process monitored by bounding hyper-rectangle $\mathcal{B}_m = \prod_{d=1}^D [L_{m,d}, R_{m,d}]$ with per-channel Bonferroni tail scaling $\alpha_{\text{dim}} = \frac{\alpha_{\text{sys}}}{D}$, the overall system false alarm probability under in-control conditions is upper-bounded by $\alpha_{\text{sys}}$:
\begin{equation}
P\left( \exists d \in \{1,\dots,D\}: x_d \notin [L_d, R_d] \mid \text{In-Control} \right) \le \alpha_{\text{sys}}
\end{equation}
\end{theorem}
\begin{proof}
Let $E_d$ denote the event that the $d^{\text{th}}$ sensor dimension exceeds its marginal control boundaries under in-control conditions:
\begin{equation}
E_d = \{ x_d \notin [L_d, R_d] \mid \text{In-Control} \}, \quad \text{for } d = 1, \dots, D
\end{equation}
By definition of the per-channel control limits $[L_d, R_d]$ using Bonferroni-scaled quantiles $\alpha_{\text{dim}} = \frac{\alpha_{\text{sys}}}{D}$, the marginal violation probability per channel is:
\begin{align}
P(E_d) &= P(x_d < L_d) + P(x_d > R_d) \notag \\
&= \frac{\alpha_{\text{sys}}}{2D} + \frac{\alpha_{\text{sys}}}{2D} = \frac{\alpha_{\text{sys}}}{D}
\end{align}
A system-wide false alarm occurs if at least one sensor dimension experiences a boundary violation, corresponding to the union event $\bigcup_{d=1}^D E_d$. By Boole's inequality (the union bound), the joint probability of the union of arbitrary events is strictly upper-bounded by the sum of their marginal probabilities, regardless of the underlying cross-channel correlation structure $\boldsymbol{\Sigma}$:
\begin{align}
P\left( \exists d \in \{1,\dots,D\}: x_d \notin [L_d, R_d] \mid \text{In-Control} \right) &= P\left( \bigcup_{d=1}^D E_d \right) \notag \\
&\le \sum_{d=1}^D P(E_d) \notag \\
&= \sum_{d=1}^D \frac{\alpha_{\text{sys}}}{D} = D \cdot \frac{\alpha_{\text{sys}}}{D} = \alpha_{\text{sys}}
\end{align}
Furthermore, if the $D$ channels are mutually independent, the exact Family-Wise Error Rate simplifies via Šidák's identity:
\begin{equation}
\text{FWER}_{\text{indep}} = 1 - \prod_{d=1}^D \left(1 - P(E_d)\right) = 1 - \left(1 - \frac{\alpha_{\text{sys}}}{D}\right)^D < \alpha_{\text{sys}}
\end{equation}
Since $\text{FWER}_{\text{indep}} < \sum_{d=1}^D P(E_d) = \alpha_{\text{sys}}$, Boole's inequality guarantees that $\alpha_{\text{sys}}$ remains a strict upper bound under arbitrary sensor cross-correlations, completing the proof.
\end{proof}

\subsection{Asymptotic Complexity Analysis}
\begin{itemize}
    \item \textbf{Space Complexity:} Strictly $O(D)$ bounded RAM ($3.23\text{ KB}$ for $D=34$), independent of stream length $N$.
    \item \textbf{Time Complexity:} Amortized $O(1)$ per sample ($< 45\text{ ms}$ per streaming batch).
\end{itemize}


% =========================================================================
% SECTION 5: TWO-TIER EMPIRICAL BENCHMARK EXPERIMENTS
% =========================================================================
\section{Two-Tier Empirical Benchmark Experiments}\label{sec5_experiments}
This section presents empirical evaluation results across a two-tier benchmark suite.

\subsection{Experimental Setup \& Evaluation Metrics}
Formulation of Coverage Rate (\%), Sample-level FAR (\%), Chunk-level FAR (\%), $\text{ARL}_0$, $\text{ARL}_1$ (Detection Delay), Peak Memory Footprint (KB), and Avg Streaming Latency (ms).

\subsection{Tier 1: Univariate Range Approximation Benchmarks}
Performance on 10 simulated 1D distributions ($F, Wald, Normal, \chi^2, Weibull$) and 4 real-world economic datasets.

\subsection{Tier 2: Real-World Industrial Telemetry Benchmarks}
Benchmark results across AI4I 2020 Predictive Maintenance ($N=10k, D=5$), MetroPT-3 Air Compressor ($N=1.51M, D=7$), Large Industrial Pump ($N=20k, D=5$), and Water Pump Sensor ($N=220k, D=10$).

\subsection{Tier 2 Stress Benchmark: Tennessee Eastman Process (TEP Multi-Mode)}
Comparative performance across TEP Modes 1 (Nominal), 3 (Feed Skewness 90/10), 4 (Max Production Throughput), and 5 (Combined Extreme Stress) for $D=34$ sensor channels.

\subsection{Hyperparameter Sensitivity Analysis}
Evaluation of chunk alarm threshold $C_{\text{thresh}} \in \{5, 10, 15\}$ on TEP Mode 1 stream.


% =========================================================================
% SECTION 6: DISCUSSION & DATA TAXONOMY INSIGHTS
% =========================================================================
\section{Discussion and Data Taxonomy Insights}\label{sec6_discussion}
Comparative discussion of non-Gaussian tail fitting, memory explosion prevention, EWMA lag suppression, and edge IoT deployment guidelines.


% =========================================================================
% SECTION 7: CONCLUSION & FUTURE WORK
% =========================================================================
\section{Conclusion}\label{sec7_conclusion}
Summary of key scientific findings and roadmap for Paper 2 (concept drift adaptation).

\bibliographystyle{cas-model2-names}
\bibliography{cas-refs}

\end{document}
